% Part: applied-modal-logics
% Chapter: temporal-logic
% Section: introduction

\documentclass[../../../include/open-logic-section]{subfiles}

\begin{document}

\olfileid{aml}{tl}{int}

\olsection{مقدمة}

يتناول المنطق الزمني دعاوى عما سيكون واقعًا أو كان واقعًا. ويُنسب إلى آرثر
برايور تأسيس المنطق الزمني، الذي سماه \emph{منطق الأزمنة}. وسيتبع تناولنا
للمنطق الزمني هنا إلى حد بعيد معالجة برايور الجهوية الأصلية، التي تُدخل
المؤثرات الزمنية في الإطار الأساسي للمنطق القضي، وهو إطار يعامل الدعاوى في
العموم بوصفها غير مقيدة بزمن.

فمثلًا، قد أتحدث في المنطق القضي عن كلبة اسمها بيزي، تجلس أحيانًا ولا تجلس
أحيانًا أخرى، شأن الكلاب عادةً. وسيكون من المتناقض في المنطق الكلاسيكي أن
نزعم أن بيزي جالسة وأنها ليست جالسة أيضًا. لكن من الواضح أن القضيتين قد
تصدقان، وإن لم تصدقا في الوقت نفسه؛ وإضافة المؤثرات الزمنية إلى اللغة تتيح لنا
التعبير عن هذه الدعوى بسهولة نسبية. كما تتيح لنا إضافة المؤثرات الزمنية
تفسير صلاحية استدلالات من قبيل الانتقال من «ستحصل بيزي على مكافأة أو كرة»
إلى «ستحصل بيزي على مكافأة أو ستحصل بيزي على كرة».

غير أن المنطق الزمني يثير قضايا فلسفية كثيرة قد تدفعنا إلى اعتماد إطار له
دون آخر. فالحادث المستقبلي، مثلًا، عبارة عن المستقبل ليست ضرورية ولا
مستحيلة. فإذا قلنا «سيذهب ريتشارد إلى متجر البقالة غدًا»، عبّرنا عن دعوى
تتعلق بشيء لم يقع بعد، وقيمة صدقها موضع نزاع. بل إن إمكان \emph{إسناد}
قيمة صدق إلى هذه الدعوى أصلًا موضع نزاع. فإذا كنا حتميين صارمين، أمكن أن
نطمئن إلى فكرة أن هذه الجملة إما صادقة وإما كاذبة بالفعل، حتى قبل الوقت
المفترض لوقوع الحدث المعني؛ غير أننا قد لا نعرف قيمة صدقها بعد. وفي المقابل،
قد نؤمن بمستقبل مفتوح حقًا، لا تتحدد فيه قيم صدق الحوادث المستقبلية.

ويتبين أن كثيرًا من هذه الالتزامات المتعلقة ببنية الزمن وطبيعته مضمّن في
اختياراتنا لنماذج المنطق الزمني وأطره. فقد نسأل، مثلًا، هل ينبغي أن نبني
نماذج يكون الزمن فيها خطيًا أو متفرعًا أو حتى دائريًا. وقد يلزم أن نقرر هل
لنماذجنا الزمنية نقطتا بداية ونهاية، وهل نمثل الزمن بلحظات منفصلة أم بمتصل.

\end{document}
